\subsection{应用：积分收敛的对数判别法}

通过适当地选取函数 g，可以从 V, p.~228 的命题 1 及其推论 1 导出一些判别法，用以判定积分 \(\int_{a}^{+\infty} f(t) dt\) （函数 \(f \geqslant 0\) 的）收敛或为无穷：只需为 g 选取一个原函数为已知的函数即可。特别地，由于 \(x^{\mu}\) 以 \(\frac{x^{\mu+1}}{\mu+1}\) 为原函数（当 \(\mu \neq -1\) 时），而以 \(\log x\) 为原函数（当 \(\mu = -1\) 时），我们有下列判别法：

命题 2（``0 阶对数判别法''）。设 f 是一个 \(\geqslant 0\) 于区间 \([a, +\infty[\) 上的规则函数；若有 \(f(x) \preccurlyeq x^{\mu}\) （对某个 \(\mu < -1\) ），则积分 \(\int_{a}^{+\infty} f(t) dt\) 收敛；若 \(f(x) \succcurlyeq x^{\mu}\) （对某个 \(\mu \geqslant -1\) ），则积分 \(\int_{a}^{+\infty} f(t) dt\) 为无穷。

当 \(1 / x^{1 + \alpha} \ll f(x) \ll 1 / x\) （对每个指数 \(\alpha > 0\) ）时，这个判别法不是决定性的，例如当 \(f(x) = 1 / x(\log x)^{\mu}\) （ \(\mu > 0\) ）时即是如此。但在后一情形， \(f\) 以 \(\frac{1}{1 - \mu} (\log x)^{1 - \mu}\) 为原函数（若 \(\mu \neq 1\) ），以 \(\log \log x\) 为原函数（当 \(\mu = 1\) 时）。于是：

命题 3（``1 阶对数判别法''）。设 f 是一个 \(\geqslant 0\) 于区间 \([a, +\infty[\) 上的规则函数；若 \(f(x) \preccurlyeq 1/x(\log x)^{\mu}\) （对某个 \(\mu > 1\) ），则积分 \(\int_{a}^{+\infty} f(t) dt\) 收敛；若 \(f(x) \succcurlyeq 1/x(\log x)^{\mu}\) （对某个 \(\mu \leqslant 1\) ），则积分 \(\int_{a}^{+\infty} f(t) dt\) 为无穷。

一般地，对每个整数 \(n \geqslant 0\) ，用 \(l_{n}(x)\) 表示由关系 \(l_{0}(x) = x\) 与 \(l_{n}(x) = \log(l_{n-1}(x))\) （对 \(n \geqslant 1\) ）归纳地定义的函数（对充分大的 x）；称 \(l_{n}(x)\) 为 x 的 \(n^{th}\) 迭对数（cf.~附录）。可以立即验证， \(\frac{1}{1 - \mu}\left(l_{n}(x)\right)^{1 - \mu}\) 是

\[
\frac {1}{x l _ {1} (x) l _ {2} (x) \dots l _ {n - 1} (x) \left(l _ {n} (x)\right) ^ {\mu}}
\]

的原函数（当 \(\mu \neq 1\) 时），而 \(l_{n+1}(x)\) 是 \(\frac{1}{xl_1(x)l_2(x)\ldots l_{n-1}(x)l_n(x)}\) 的原函数。因此：

命题 4（``n 阶对数判别法''）。设 f 是一个 \(\geqslant0\) 于区间 \([a,+\infty[\) 上的规则函数；若对某个 \(\mu>1\) 有 \(f(x)\preccurlyeq\frac{1}{xl_{1}(x)l_{2}(x)\ldots l_{n-1}(x)(l_{n}(x))^{\mu}}\) ，则积分 \(\int_{a}^{+\infty}f(t)dt\) 收敛；若 \(f(x)\succcurlyeq\frac{1}{xl_{1}(x)l_{2}(x)\ldots l_{n-1}(x)(l_{n}(x))^{\mu}}\) （对某个 \(\mu\leqslant1\) ），则积分 \(\int_{a}^{+\infty}f(t)dt\) 为无穷。

这样，每个对数判别法适用于那些使更低阶判别法不是决定性的函数（cf.~V, p.~264, exerc. 5 b) 与 V, p.~265, exerc. 8）。

由于它的用处，我们对积分 \(\int_{\alpha}^{a}f(t)dt\) 翻译 0 阶判别法，其中 \(f\) 是规则函数且 \(\geqslant 0\) 于非紧区间 \([\alpha ,a]\) 上：

命题 5（``0 阶对数判别法''）。若在 \(\alpha\) 的一个邻域上有 \(f(x) \preccurlyeq 1/(x - \alpha)^{\mu}\) （对某个 \(\mu < 1\) ），则积分 \(\int_{\alpha}^{a} f(t) dt\) 收敛；若 \(f(x) \succcurlyeq 1/(x - \alpha)^{\mu}\) （对某个 \(\mu \geqslant 1\) ），则积分 \(\int_{\alpha}^{a} f(t) dt\) 为无穷。

我们把 \(n\) 阶对数判别法的类似翻译留给读者。

如果会求 f 相对于一个包含这些判别法中所出现函数的比较尺度的主部，对数判别法的应用便是直接的：若 \(f_{1}\) 是主部，则积分 \(\int_{\alpha}^{+\infty} f(t) dt\) 与 \(\int_{\alpha}^{+\infty} f_{1}(t) dt\) 同时收敛或同时为无穷，而对数判别法可直接应用于后一积分。

例. 1) 函数 \(t^p (1 - t)^q\) 在 {]}0, 1{[} 上不是有界的（当 \(p < 0\) 或 \(q < 0\) 时）；用在点 0 和 1 的邻域上应用的 0 阶对数判别法，积分 \(\int_0^1 t^p (1 - t)^q dt\) 收敛当且仅当 \(p > -1\) 且 \(q > -1\) 。这时，这个积分称为第一类欧拉积分，记作 \(\mathbf{B}(p + 1,q + 1)\) (cf.~VII, p.~312)。

\begin{enumerate}
\def\labelenumi{\arabic{enumi})}
\setcounter{enumi}{1}
\tightlist
\item
  考察积分 \(\int_0^\infty t^{x - 1}e^{-t}dt\) 。由于 \(e^{-t}\sim 1\) 在 0 的一个邻域上，要这个积分收敛，必须有 \(x > 0\) ；这个条件也是充分的，因为在 \(+\infty\) 的一个邻域上有 \(e^{-t}\ll t^{-\mu}\) （对任何 \(\mu >0\) ）。当 \(x > 0\) 时，这个积分称为第二类欧拉积分，记作 \(\Gamma (x)\) (cf.~VII, p.~311)。
\end{enumerate}

